\section{Scientific background}
\label{sec:background}

This section connects the large questions of cosmology to the smaller
modelling problem studied in this project. We first introduce the expanding
Universe, then explain what galaxy shapes can tell us and why their intrinsic
alignments matter.

% STUDENT FIGURES: detailed panel layouts, source links, and checks are in
% Report/Section2_Figure_Guide.md. Each slot below becomes an image when its
% named PDF is placed in Report/figures/. The boxes are drafting aids.
% This macro is used only in Section 2; later report figures are unchanged.
\newcommand{\backgroundfigure}[4]{%
  \begin{figure}[htbp]
    \centering
    \IfFileExists{figures/#1}{%
      \includegraphics[width=0.96\linewidth,height=0.34\textheight,
        keepaspectratio]{#1}%
    }{%
      \fbox{\parbox[c][3.0cm][c]{0.90\linewidth}{%
        \centering\small\textbf{Student figure to be prepared}\par
        \medskip #2\par\medskip
        \texttt{\detokenize{#1}}\par
        See \texttt{Section2\_Figure\_Guide.md}.}}%
    }
    \caption{#3}
    \label{#4}
  \end{figure}%
}

\subsection{An expanding Universe and its contents}
\label{sec:cosmology-background}

Cosmology studies the Universe as a whole: its contents, history, and large
structures. The hot Big Bang model describes an early Universe that was much
hotter and denser than it is today. As space expanded, the Universe cooled.
The Big Bang was not an explosion from one place into empty space; it
describes the expansion of space itself. A key piece of evidence is the
\emph{cosmic microwave background} (CMB), light that began travelling freely
when the Universe was about 380,000 years old. The Planck satellite mapped
its tiny temperature variations, which carry information about conditions
in the early Universe \citep{2020A&A...641A...1P}.

Expansion also stretches travelling light to longer wavelengths. We describe
this \emph{cosmological redshift} by
\[
  1+z=\frac{\lambda_{\rm observed}}{\lambda_{\rm emitted}}
      =\frac{1}{a},
\]
where $a$ is the scale factor, which tracks cosmic expansion and equals
one today. Thus $z=1$ means that the wavelength has doubled since emission.
Larger cosmological redshifts generally let us study earlier times. They
are not distances by themselves: converting redshift to distance requires
an expansion model.

The standard model, called $\Lambda$CDM, includes a cosmological constant
($\Lambda$) and cold dark matter (CDM). It brings together evidence from the
CMB, galaxy distributions, and other observations. In this model, about
$5\%$ of today's total energy density is ordinary matter, $27\%$ is dark
matter, and $68\%$ is dark energy; these are approximate, model-dependent
fractions \citep{2020A&A...641A...6P}. Ordinary matter includes atoms in stars
and gas. Dark matter is inferred through gravity, and its particle identity
remains unknown. ``Cold'' means that it moved slowly enough in the early
Universe to allow small structures to grow. The symbol $\Lambda$ represents
a cosmological constant, whose energy density remains constant as space
expands and can drive accelerated expansion. Whether dark energy is exactly
this constant, and why it exists, remain open questions.
Figure~\ref{fig:background-planck} connects the early-Universe evidence to
this present-day description.

% FIGURE BRIEF 2.1: Planck CMB map + present-day composition, with a short
% expansion/redshift strip. Preserve the map colour scale and ESA credit.
\backgroundfigure{background_planck_cosmology.pdf}
  {Planck temperature map; present-day cosmic contents; expansion and redshift.}
  {Planck's all-sky map of CMB temperature anisotropies. The colors show small
   temperature differences in light released about 380,000 years after the
   Big Bang; this is not a map of present-day dark matter. Image credit:
   ESA/Planck Collaboration \citep{2020A&A...641A...1P}.}
  {fig:background-planck}

\subsection{Is dark energy constant?}
\label{sec:dark-energy-background}

To test the expansion history, astronomers need ways to compare distances
at different times. The Dark Energy Spectroscopic Instrument (DESI) measures
galaxy and quasar redshifts over a large volume. Its measurements reveal
\emph{baryon acoustic oscillations} (BAO): a preferred separation in the
distribution of matter, left by sound waves in the early Universe. This
provides a statistical ruler whose apparent size changes with distance.

The second DESI data release, DR2, gives precise BAO measurements
\citep{2025PhRvD.112h3515A}. Two useful quantities are $\Omega_{\rm m}$, the
present matter density divided by the critical density for a spatially flat
Universe, and $H_0r_d$. Here $H_0$ is today's expansion rate and $r_d$ is the
comoving length of the BAO ruler, with the overall expansion factored out.
BAO constrains $H_0r_d$ more directly than $H_0$ alone.
Figure~\ref{fig:background-desi} shows the
allowed combinations and a test of changing dark energy.

Dark energy's \emph{equation-of-state parameter}, $w$, is its pressure
divided by its energy density. A cosmological constant has $w=-1$. A common
test allows
\[
  w(a)=w_0+w_a(1-a),
\]
where $w_0$ is the value today and $w_a$ controls how it changes. The constant
case is $(w_0,w_a)=(-1,0)$. Type Ia supernovae provide another distance
indicator: after correcting for differences between them, their brightnesses
can be compared. In the DR2 BAO analysis, combining DESI with CMB
measurements and different supernova samples gives a preference for evolution
at about $2.8$--$4.2\sigma$, depending on the combination
\citep{2025PhRvD.112h3515A}. Here $\sigma$ expresses statistical significance
on a standard-deviation scale. This indication depends on the data and model
assumptions, so it is not a confirmed discovery. A 2026 DESI analysis using
absorption by intergalactic hydrogen gives results consistent with $\Lambda$CDM,
showing that the wider picture is still developing
\citep{2026arXiv260727410D}.

Confirmed evolution would challenge the cosmological-constant explanation
and have major implications for fundamental physics. Weak lensing offers
complementary evidence through the growth and distribution of matter, with
different sources of measurement error. Such a check must account for shared
data and assumptions; a combined analysis is not automatically independent.

% FIGURE BRIEF 2.2: DESI DR2 BAO paper, Figure 8 LEFT panel and Figure 11.
% Use released contours or faithfully attributed panels, never invented ovals.
\backgroundfigure{background_desi_dr2.pdf}
  {Two panels: $\Omega_{\rm m}$--$H_0r_d$ and $w_0$--$w_a$ constraints.}
  {Two tests using DESI DR2 BAO measurements. Left: allowed values of
   $\Omega_{\rm m}$ and $H_0r_d$ assuming $\Lambda$CDM. Right: allowed
   $w_0$ and $w_a$ values for the labelled DESI, CMB, and supernova
   combinations. Inner and outer contours enclose $68\%$ and $95\%$ of the
   posterior probability under each model. The point $(-1,0)$ marks a
   cosmological constant. Source: \citet{2025PhRvD.112h3515A},
   Figures 8 (left panel) and 11.}
  {fig:background-desi}

\subsection{Weak lensing and the distribution of matter}
\label{sec:lensing-background}

Gravity changes the paths of light rays. A large mass can produce obvious
arcs or multiple images of a background galaxy, called strong lensing. In
weak lensing, the changes are much smaller. \emph{Convergence}, written
$\kappa$, describes isotropic focusing: a small circular image changes
size but stays circular. \emph{Shear}, written $\gamma$, describes stretching
in one direction and compression in the perpendicular direction. A circle
then becomes an ellipse, as illustrated in Figure~\ref{fig:background-lensing}
\citep{2001PhR...340..291B}.

The matter doing the lensing is not spread uniformly. It forms clusters,
filaments, and emptier regions. We describe this \emph{density field} by its
density contrast,
\[
  \delta(\boldsymbol{x},z)=
  \frac{\rho(\boldsymbol{x},z)-\bar\rho(z)}{\bar\rho(z)},
\]
where $\rho$ is the matter density at position $\boldsymbol{x}$ and
$\bar\rho$ is its mean at that time. Positive $\delta$ denotes an overdense
region. Instead of tracking every structure, we can describe the typical
strength of fluctuations on different scales with the \emph{matter power
spectrum}, $P_\delta(k,z)$. The wavenumber $k=2\pi/\lambda_{\rm spatial}$
labels a spatial wavelength: small $k$ means large structures, and large
$k$ means small structures. We use comoving lengths, which factor out the
overall expansion. A megaparsec (Mpc) is about 3.26 million light-years,
so $k$ can be expressed in $\mathrm{Mpc}^{-1}$. The power spectrum summarizes
how densities at pairs of positions are related, called two-point statistics;
it does not specify the full arrangement of matter
\citep{2001PhR...340..291B}.

Astronomers estimate galaxy shapes, correct for image blurring and measurement
bias, and estimate galaxy redshifts. They then measure how pairs of shapes
are correlated at different angular separations and in different redshift
groups. This statistical lensing signal is called \emph{cosmic shear}.
Its prediction combines the matter power spectrum with the geometry between
the observer, intervening matter, and source galaxies. Matter over a broad
range of distances contributes to each measurement. Lensing therefore tests
both the growth of structure and the expansion history, but does not directly
give $P_\delta$ at a single scale and redshift.

% FIGURE BRIEF 2.3: source--matter--observer geometry above; identical source
% circles under convergence and shear below; add a small scale-to-k inset.
\backgroundfigure{background_lensing_density.pdf}
  {Light-path geometry; convergence and shear; large and small spatial scales.}
  {How weak lensing connects images to matter. Intervening matter deflects
   light from distant galaxies. Convergence changes apparent size, while
   shear changes shape; distortions are exaggerated for clarity. A scale
   inset connects broad and fine density variations to small and large
   wavenumbers. This is a schematic explanation, not measured galaxy data
   \citep{2001PhR...340..291B}. This AI-generated original schematic was
   produced with GPT-5.6 Sol.}
  {fig:background-lensing}

\subsection{What galaxy surveys tell us, and what comes next}
\label{sec:surveys-background}

Three major imaging surveys are the Kilo-Degree Survey (KiDS), the Dark
Energy Survey (DES), and the Hyper Suprime-Cam survey (HSC). Often called
Stage III surveys, they use different combinations of area, depth, and image
quality. A useful quantity for comparing their cosmic-shear results is
\[
  S_8=\sigma_8\sqrt{\Omega_{\rm m}/0.3}.
\]
Here $\sigma_8$ measures the root-mean-square density contrast predicted
today by linear growth theory, after averaging within spheres of radius
$8h^{-1}\,\mathrm{Mpc}$, with
$h=H_0/(100\,\mathrm{km\,s^{-1}\,Mpc^{-1}})$. Larger $\sigma_8$ means stronger
matter fluctuations. Cosmic shear constrains the combination $S_8$
particularly well.

Some earlier shear analyses favoured a lower $S_8$ than the value predicted
from the CMB within $\Lambda$CDM. This difference is known as the
\emph{$S_8$ tension}. The KiDS-1000 analysis found a difference of about
$3\sigma$ \citep{2021A&A...645A.104A}; its later reanalysis reduced this to
about $2.3\sigma$ \citep{2023A&A...679A.133L}. KiDS-Legacy is now consistent
with Planck, with $S_8=0.815^{+0.016}_{-0.021}$. The shift mainly reflects
improved redshift calibration, together with new survey area and improved
image reduction \citep{2025A&A...703A.158W}.

The situation is not identical in every survey. DES Year 6 reports
$S_8=0.798^{+0.014}_{-0.015}$ or $0.783^{+0.019}_{-0.015}$ for two different
alignment models, discussed in Section~\ref{sec:ia-models-background}
\citep{2026arXiv260210065D}. These are alternative analyses of the same
observations, not two independent measurements. HSC Year 3 finds
$S_8=0.769^{+0.031}_{-0.034}$ from shape correlation functions
\citep{2023PhRvD.108l3518L}. Differences of roughly $2\sigma$ remain in
some comparisons. Figure~\ref{fig:background-surveys} compares successive
analyses, including DES Year 3 and HSC Year 1
\citep{2022PhRvD.105b3514A,2020PASJ...72...16H}. The lesson is that the
tension has eased in some analyses; it has not simply disappeared everywhere.

The next generation, often called Stage IV, includes Euclid, the Vera C.
Rubin Observatory's Legacy Survey of Space and Time, and the Nancy Grace
Roman Space Telescope. Their programmes combine wide sky coverage, repeated
imaging, and sharp space-based images in different ways
\citep{2025A&A...697A...1E,2019ApJ...873..111I,2021MNRAS.507.1746E}.
Smaller statistical errors make systematic errors more important. These
include shape and redshift calibration, the effect of gas and stars on the
matter distribution, and intrinsic alignment. Collecting more galaxies alone
does not remove these uncertainties.

% FIGURE BRIEF 2.4: Stage III measured contours, NOT Stage IV detections.
% Compare like statistics and label model/analysis changes. Keep prospects
% separate from measurements. See the guide for the KiDS-1000-v2 distinction.
\backgroundfigure{background_s8_surveys.pdf}
  {Measured Stage III $S_8$ contours across releases; separate Stage IV prospects.}
  {Published $\Omega_{\rm m}$--$S_8$ constraints from KiDS-Legacy, DES Year 3,
   HSC Year 3, and Planck-Legacy. Inner and outer contours show the reported
   probability regions. This comparison predates the DES Year 6 result
   discussed in the text. Adapted from Figure 14 of
   \citet{2025A&A...703A.158W} under CC BY 4.0.}
  {fig:background-surveys}

\subsection{Intrinsic alignment: how environments affect galaxy shapes}
\label{sec:ia-background}

Galaxies do not form in isolation. Gravity pulls differently across a galaxy
and its surroundings; these differences are called a \emph{tidal field}.
They can influence galaxy shapes and orientations. Tidal forces can also
exert torques that affect rotation. These mechanisms help explain why
intrinsic alignments depend on galaxy population, environment, scale, and
time. Observations and simulations support parts of this picture, but no
single model describes every population and scale equally well
\citep{2025A&ARv..33....5C}.

An image's \emph{ellipticity} records its elongation and orientation. To
understand the contributions, we can write schematically
\[
 \underbrace{e^{\rm obs}}_{\text{measured shape}}
 \simeq
 \underbrace{e^{\rm I}}_{\text{intrinsic shape}}
 +\underbrace{\gamma}_{\text{gravitational shear}}
 +\underbrace{\eta}_{\text{measurement noise}}.
\]
Each shape or shear has two components describing magnitude and orientation.
This is an illustration for weak distortions after accounting for the shape
estimator's response, not an exact rule for adding arbitrary ellipses.
The more precise lensing quantity is the reduced shear
$g=\gamma/(1-\kappa)$, which approaches $\gamma$ when convergence is small
\citep{2001PhR...340..291B}.

What a telescope records is a pixelated, point-spread-function-convolved
image, not a galaxy's unlensed shape or its shear. After image calibration,
the catalogue contains an estimate of each galaxy's apparent ellipticity,
sky position, and redshift information. The intrinsic ellipticity of one
galaxy is normally much larger than the weak shear acting on it, so a single
image cannot provide a useful shear measurement. The observable lensing
signal is instead a statistical correlation of apparent ellipticities over
many galaxy pairs, binned by angular separation and estimated redshift.

Ignoring terms whose expectation is zero, expanding the product
$(G+I)_i(G+I)_j$ explains why several correlations appear. Suppressing the
angular and ellipticity-component labels, the measured two-point signal is
\[
  \langle e_i^{\rm obs}e_j^{\rm obs}\rangle_{\rm signal}
  \simeq GG_{ij}+GI_{ij}+IG_{ij}+II_{ij}.
\]
The brackets mean an average over pairs, and $i,j$ label redshift groups.
The terms have different physical origins:
\begin{itemize}
  \item \textbf{GG} correlates the gravitational shears of two source
  galaxies. Their light rays pass through correlated foreground structure,
  making GG the desired cosmic-shear contribution.
  \item \textbf{II} correlates two intrinsic orientations. Galaxies that
  formed in the same or correlated tidal environment can point in related
  directions even before lensing; this contribution is usually most relevant
  for pairs that are physically close.
  \item \textbf{GI and IG} correlate one intrinsic orientation with the
  shear of the other source. A foreground galaxy can align with surrounding
  matter that also lenses a more distant galaxy. The ordering identifies
  which redshift bin supplies $G$ and which supplies $I$. The two orderings
  are commonly grouped under the name GI.
\end{itemize}
The survey measures only their sum, not four separately labelled signals.
Their different redshift and angular dependences help a statistical model
distinguish them, but the separation is model-dependent. Random intrinsic
shapes and uncorrelated measurement noise mainly add variance (``shape
noise''); correlated measurement errors can instead bias the result.
Figure~\ref{fig:ia-clue} illustrates the three physical connections
\citep{2025A&ARv..33....5C}.

% FIGURE BRIEF 2.5: draw GG, II, and GI as three different geometries; put
% the schematic ellipticity sum above them. Retain the existing figure label.
\backgroundfigure{background_intrinsic_alignment.pdf}
  {Intrinsic shape plus shear; separate GG, II, and GI pair diagrams.}
  {Why observed shapes contain more than lensing. The shape sum is schematic.
   GG links two lensed images, II links two intrinsic orientations, and GI
   links a foreground intrinsic orientation to the shear of a background
   source. Here GI groups the two cross terms when both orderings are
   relevant. The diagram illustrates mechanisms, not their relative strengths
   \citep{2025A&ARv..33....5C}. This AI-generated original schematic was
   produced with GPT-5.6 Sol.}
  {fig:ia-clue}

\subsection{The modelling tradeoff and our compression question}
\label{sec:ia-models-background}

Several models have been used to describe intrinsic alignment. The
\emph{linear alignment} (LA) picture treats a galaxy's intrinsic shape as a
leading-order response to the large-scale tidal field. The widely used
\emph{nonlinear alignment} (NLA) prescription keeps this basic response
structure but substitutes a nonlinear matter power spectrum when predicting
the correlations. Typical survey implementations then describe the
amplitude and redshift evolution with only a few nuisance parameters
\citep{2007NJPh....9..444B}. NLA is useful and economical, but the word
``nonlinear'' does not make it a complete model of nonlinear galaxy
formation or guarantee accuracy on every scale.

The \emph{tidal alignment and tidal torquing} (TATT) model includes both a
leading tidal-alignment response and additional responses, including terms
quadratic in the tidal field and effects associated with how galaxies sample
the density field \citep{2019PhRvD.100j3506B}. It can therefore represent
behaviour that NLA cannot, but it normally introduces more adjustable
parameters. Neither model is universally best: suitability depends on the
galaxy sample, physical scales, statistical precision, and priors used in the
analysis.

This creates the central modelling difficulty. If the alignment model is too
restricted, an unmodelled IA contribution can be mistaken for a change in
matter clustering and bias parameters such as $S_8$. If the model is very
flexible, several combinations of IA, redshift uncertainty, and cosmology can
predict similar projected shape correlations. These \emph{parameter
degeneracies} weaken constraints, make results more dependent on priors, and
increase the cost of exploring the parameter space. Projection along the
line of sight further hides differences between scale and redshift
dependences. The DES Year 6 comparison illustrates that changing the
alignment prescription can shift the inferred $S_8$
\citep{2026arXiv260210065D}.

This project is motivated by the possibility that a flexible response may
have many sampled parameters but vary mainly along a much smaller number of
effective directions. As shown in
Figure~\ref{fig:background-model-tradeoff}, 13 shape parameters generate each
positive response surface $\mathcal A_\Theta(k,z)$ on a $31\times101$ grid.
We test whether PCA or an autoencoder can reconstruct this family using only
two coordinates. If successful, such a representation could eventually
reduce the effective number of nuisance coordinates that must be explored
while retaining more scale--redshift flexibility than a very restricted
model.

The scope of that claim is deliberately limited. The synthetic family is
NLA-inspired; it is not a TATT emulator, the matter power spectrum, or a
complete physical IA model. Two accurate reconstruction coordinates would
not prove that IA has only two physical parameters, remove degeneracies, or
guarantee unbiased cosmology. A practical analysis would still need to
propagate reconstruction errors into lensing predictions, define justified
probabilities for the compressed coordinates, and test inference on realistic
survey data. Section~\ref{sec:model} defines the controlled dataset used for
this first compression test.

% FIGURE BRIEF 2.6: two balanced model branches converge on the research
% question. Mark reconstruction as tested and cosmological inference as future.
\clearpage
\begingroup
\makeatletter
\setlength{\@fptop}{0pt}
\makeatother
\backgroundfigure{background_model_tradeoff.pdf}
  {Simple and flexible models; benefits and limits; our bounded research question.}
  {The modelling tradeoff motivating this project. Restricted models are
   easier to constrain but allow fewer responses; flexible models allow
   more responses but can introduce poorly constrained directions and
   additional computation. Our experiment tests compact reconstruction of
   a chosen response family. It does not yet demonstrate faster or unbiased
   cosmological inference. NLA and TATT provide the physical context
   \citep{2007NJPh....9..444B,2019PhRvD.100j3506B}. This AI-generated
   original schematic was produced with GPT-5.6 Sol.}
  {fig:background-model-tradeoff}

% Keep the background's figures before Section 3.
\clearpage
\endgroup
